\section{Warsztat EDC}

% SLIDE ID: sec03-goal
\begin{frame}{Cel warsztatu}
  \SlideID{sec03-goal}
  \begin{itemize}
    \item wejść do systemu jako `CRC / site user`,
    \item wprowadzić dane pacjenta do dwóch wizyt,
    \item zobaczyć, jak system reaguje na błędy, braki i query.
  \end{itemize}
\end{frame}

% SLIDE ID: sec03-login
\begin{frame}{Instrukcja logowania}
  \SlideID{sec03-login}
  \begin{enumerate}
    \item otwórz środowisko szkoleniowe EDC,
    \item zaloguj się danymi przekazanymi przez prowadzącego,
    \item wybierz ośrodek i listę pacjentów,
    \item otwórz przypisany rekord pacjenta szkoleniowego.
  \end{enumerate}
\end{frame}

% SLIDE ID: sec03-patient
\begin{frame}{Scenariusz pacjenta}
  \SlideID{sec03-patient}
  \begin{itemize}
    \item inicjały: `A.K.`
    \item wiek: `54 lata`
    \item płeć: kobieta
    \item wizyta `Screening`: kwalifikacja i badania wstępne
    \item wizyta `Baseline`: potwierdzenie włączenia i dane startowe
  \end{itemize}
\end{frame}

% SLIDE ID: sec03-data-entry
\begin{frame}{Dane do wprowadzenia}
  \SlideID{sec03-data-entry}
  \small
  \begin{tabularx}{\linewidth}{>{\raggedright\arraybackslash}p{0.20\linewidth} >{\raggedright\arraybackslash}p{0.33\linewidth} >{\raggedright\arraybackslash}X}
    \toprule
    Wizyta & Pole & Wartość \\
    \midrule
    Screening & Data zgody & 2026-05-12 \\
    Screening & Data wizyty & 2026-05-10 \\
    Screening & Hemoglobina & 13.2 g/dL \\
    Screening & ALT & 96 U/L \\
    Screening & Kreatynina & brak \\
    Baseline & Data wizyty & 2026-05-18 \\
    Baseline & Randomizacja & Tak \\
    Baseline & Ciśnienie tętnicze & 128/78 mmHg \\
    \bottomrule
  \end{tabularx}
\end{frame}

% SLIDE ID: sec03-built-in-issues
\begin{frame}{Celowo wbudowane problemy / błędy}
  \SlideID{sec03-built-in-issues}
  \begin{itemize}
    \item jeden wynik poza zakresem: `ALT = 96 U/L`,
    \item jeden brak danych: brak kreatyniny,
    \item jeden błąd logiczny: data wizyty `Screening` przed datą zgody.
  \end{itemize}
\end{frame}

% SLIDE ID: sec03-participant-task
\begin{frame}{Co uczestnik ma zrobić}
  \SlideID{sec03-participant-task}
  \begin{enumerate}
    \item odnaleźć pacjenta i otworzyć właściwe wizyty,
    \item wprowadzić wszystkie dostępne dane,
    \item zapisać formularze i sprawdzić komunikaty systemu,
    \item wskazać, które problemy wymagają poprawy lub wyjaśnienia.
  \end{enumerate}
\end{frame}

% SLIDE ID: sec03-query-response
\begin{frame}{Jak reagować na query}
  \SlideID{sec03-query-response}
  \begin{itemize}
    \item przeczytaj dokładnie, czego dotyczy pytanie,
    \item wróć do danych źródłowych lub scenariusza pacjenta,
    \item popraw wpis albo dodaj sensowny komentarz,
    \item nie zamykaj query „na ślepo”, bez zrozumienia problemu.
  \end{itemize}
\end{frame}

% SLIDE ID: sec03-trainer-notes
\begin{frame}{Na co zwrócić uwagę podczas ćwiczenia}
  \SlideID{sec03-trainer-notes}
  \begin{itemize}
    \item czy uczestnik trafia do właściwej wizyty i formularza,
    \item czy rozumie różnicę między błędem, brakiem danych i query,
    \item czy potrafi logicznie uzasadnić poprawkę,
    \item czy korzysta z systemu jak użytkownik ośrodka, a nie jak data manager.
  \end{itemize}
\end{frame}

% SLIDE ID: sec03-summary
\begin{frame}{Podsumowanie: czego nauczył nas warsztat}
  \SlideID{sec03-summary}
  \begin{itemize}
    \item poprawne wprowadzanie danych wymaga rozumienia struktury badania,
    \item walidacje pomagają, ale nie zastępują myślenia użytkownika,
    \item jakość danych zależy od uważności już na poziomie ośrodka.
  \end{itemize}
\end{frame}
